
\documentclass{article}
\usepackage{colortbl}
\usepackage{makecell}
\usepackage{multirow}
\usepackage{supertabular}

\begin{document}

\newcounter{utterance}

\twocolumn

{ \footnotesize  \setcounter{utterance}{1}
\setlength{\tabcolsep}{0pt}
\begin{supertabular}{c@{$\;$}|p{.15\linewidth}@{}p{.15\linewidth}p{.15\linewidth}p{.15\linewidth}p{.15\linewidth}p{.15\linewidth}}

    \# & $\;$A & \multicolumn{4}{c}{Game Master} & $\;\:$B\\
    \hline 

    \theutterance \stepcounter{utterance}  

    & & \multicolumn{4}{p{0.6\linewidth}}{\cellcolor[rgb]{0.9,0.9,0.9}{%
	\makecell[{{p{\linewidth}}}]{% 
	  \tt {\tiny [A$\langle$GM]}  
	 Du nimmst an einem kollaborativen Verhandlungspiel Teil.\\ \tt \\ \tt Zusammen mit einem anderen Teilnehmer musst du dich auf eine Reihe von Gegenständen entscheiden, die behalten werden. Jeder von euch hat eine persönliche Verteilung über die Wichtigkeit der einzelnen Gegenstände. Jeder von euch hat eine eigene Meinung darüber, wie wichtig jeder einzelne Gegenstand ist (Gegenstandswichtigkeit). Du kennst die Wichtigkeitsverteilung des anderen Spielers nicht. Zusätzlich siehst du, wie viel Aufwand jeder Gegenstand verursacht.  \\ \tt Ihr dürft euch nur auf eine Reihe von Gegenständen einigen, wenn der Gesamtaufwand der ausgewählten Gegenstände den Maximalaufwand nicht überschreitet:\\ \tt \\ \tt Maximalaufwand = 2984\\ \tt \\ \tt Hier sind die einzelnen Aufwände der Gegenstände:\\ \tt \\ \tt Aufwand der Gegenstände = {"C76": 350, "C38": 276, "C56": 616, "A03": 737, "A07": 531, "A83": 389, "C98": 24, "C08": 125, "C62": 338, "C00": 356, "C10": 143, "C81": 117, "B38": 257, "C03": 921, "C32": 789}\\ \tt \\ \tt Hier ist deine persönliche Verteilung der Wichtigkeit der einzelnen Gegenstände:\\ \tt \\ \tt Werte der Gegenstandswichtigkeit = {"C76": 138, "C38": 583, "C56": 868, "A03": 822, "A07": 783, "A83": 65, "C98": 262, "C08": 121, "C62": 508, "C00": 780, "C10": 461, "C81": 484, "B38": 668, "C03": 389, "C32": 808}\\ \tt \\ \tt Ziel:\\ \tt \\ \tt Dein Ziel ist es, eine Reihe von Gegenständen auszuhandeln, die dir möglichst viel bringt (d. h. Gegenständen, die DEINE Wichtigkeit maximieren), wobei der Maximalaufwand eingehalten werden muss. Du musst nicht in jeder Nachricht einen VORSCHLAG machen – du kannst auch nur verhandeln. Alle Taktiken sind erlaubt!\\ \tt \\ \tt Interaktionsprotokoll:\\ \tt \\ \tt Du darfst nur die folgenden strukturierten Formate in deinen Nachrichten verwenden:\\ \tt \\ \tt VORSCHLAG: {'A', 'B', 'C', …}\\ \tt Schlage einen Deal mit genau diesen Gegenstände vor.\\ \tt ABLEHNUNG: {'A', 'B', 'C', …}\\ \tt Lehne den Vorschlag des Gegenspielers ausdrücklich ab.\\ \tt ARGUMENT: {'...'}\\ \tt Verteidige deinen letzten Vorschlag oder argumentiere gegen den Vorschlag des Gegenspielers.\\ \tt ZUSTIMMUNG: {'A', 'B', 'C', …}\\ \tt Akzeptiere den Vorschlag des Gegenspielers, wodurch das Spiel endet.\\ \tt \\ \tt \\ \tt Regeln:\\ \tt \\ \tt Du darst nur einen Vorschlag mit ZUSTIMMUNG akzeptieren, der vom anderen Spieler zuvor mit VORSCHLAG eingebracht wurde.\\ \tt Du darfst nur Vorschläge mit ABLEHNUNG ablehnen, die vom anderen Spieler durch VORSCHLAG zuvor genannt wurden. \\ \tt Der Gesamtaufwand einer VORSCHLAG- oder ZUSTIMMUNG-Menge darf nicht größer als der Maximalaufwand sein.  \\ \tt Offenbare deine versteckte Wichtigkeitsverteilung nicht.\\ \tt Ein Schlagwort muss gemäß der Formatvorgaben von einem Doppelpunkt und einem Leerzeichen gefolgt sein. Das Argument muss eine Python-Menge mit 0 oder mehr Strings sein.  \\ \tt Die Form muss also SCHLAGWORT: {...} sein.\\ \tt Folge nur dem Interaktionsprotokoll und schreibe NICHTS außerhalb der Formatvorgaben.\\ \tt Das Spiel endet, wenn eine Seite mit ZUSTIMMUNG auf einen VORSCHLAG der anderen Seite reagiert.  \\ \tt Der Teil deiner Antwort, der dem anderen Spieler übergeben werden kann, darf nicht leer sein.  \\ \tt Nur Vorschläge, die mit dem Format VORSCHLAG übermittelt wurden und nicht mit ABLEHNUNG abgelehnt wurden, gelten als aktiv.  \\ \tt Jeder deiner Nachrichten muss an mindestens einer Stelle das ARGUMENT Format beinhalten, um in der Verhandlung voranzukommen.\\ \tt  
	  } 
	   } 
	   } 
	 & \\ 
 

    \theutterance \stepcounter{utterance}  

    & \multicolumn{4}{p{0.6\linewidth}}{\cellcolor[rgb]{0.8,1,0.9}{%
	\makecell[{{p{\linewidth}}}]{% 
	  \tt {\tiny [A$\rangle$GM]}  
	 VORSCHLAG: {'C98', 'C81', 'C10', 'C00', 'C38', 'B38', 'C56', 'A07', 'C62', 'C08'}\\ \tt ARGUMENT: {} 
	  } 
	   } 
	   } 
	 & & \\ 
 

    \theutterance \stepcounter{utterance}  

    & & \multicolumn{4}{p{0.6\linewidth}}{\cellcolor[rgb]{0.9,0.9,0.9}{%
	\makecell[{{p{\linewidth}}}]{% 
	  \tt {\tiny [GM$\rangle$B]}  
	 Du nimmst an einem kollaborativen Verhandlungspiel Teil.\\ \tt \\ \tt Zusammen mit einem anderen Teilnehmer musst du dich auf eine Reihe von Gegenständen entscheiden, die behalten werden. Jeder von euch hat eine persönliche Verteilung über die Wichtigkeit der einzelnen Gegenstände. Jeder von euch hat eine eigene Meinung darüber, wie wichtig jeder einzelne Gegenstand ist (Gegenstandswichtigkeit). Du kennst die Wichtigkeitsverteilung des anderen Spielers nicht. Zusätzlich siehst du, wie viel Aufwand jeder Gegenstand verursacht.  \\ \tt Ihr dürft euch nur auf eine Reihe von Gegenständen einigen, wenn der Gesamtaufwand der ausgewählten Gegenstände den Maximalaufwand nicht überschreitet:\\ \tt \\ \tt Maximalaufwand = 2984\\ \tt \\ \tt Hier sind die einzelnen Aufwände der Gegenstände:\\ \tt \\ \tt Aufwand der Gegenstände = {"C76": 350, "C38": 276, "C56": 616, "A03": 737, "A07": 531, "A83": 389, "C98": 24, "C08": 125, "C62": 338, "C00": 356, "C10": 143, "C81": 117, "B38": 257, "C03": 921, "C32": 789}\\ \tt \\ \tt Hier ist deine persönliche Verteilung der Wichtigkeit der einzelnen Gegenstände:\\ \tt \\ \tt Werte der Gegenstandswichtigkeit = {"C56": 65, "A03": 121, "C32": 138, "A07": 262, "C00": 389, "B38": 461, "C38": 484, "C62": 508, "C81": 583, "C10": 668, "C03": 780, "C98": 783, "C76": 808, "C08": 822, "A83": 868}\\ \tt \\ \tt Ziel:\\ \tt \\ \tt Dein Ziel ist es, eine Reihe von Gegenständen auszuhandeln, die dir möglichst viel bringt (d. h. Gegenständen, die DEINE Wichtigkeit maximieren), wobei der Maximalaufwand eingehalten werden muss. Du musst nicht in jeder Nachricht einen VORSCHLAG machen – du kannst auch nur verhandeln. Alle Taktiken sind erlaubt!\\ \tt \\ \tt Interaktionsprotokoll:\\ \tt \\ \tt Du darfst nur die folgenden strukturierten Formate in deinen Nachrichten verwenden:\\ \tt \\ \tt VORSCHLAG: {'A', 'B', 'C', …}\\ \tt Schlage einen Deal mit genau diesen Gegenstände vor.\\ \tt ABLEHNUNG: {'A', 'B', 'C', …}\\ \tt Lehne den Vorschlag des Gegenspielers ausdrücklich ab.\\ \tt ARGUMENT: {'...'}\\ \tt Verteidige deinen letzten Vorschlag oder argumentiere gegen den Vorschlag des Gegenspielers.\\ \tt ZUSTIMMUNG: {'A', 'B', 'C', …}\\ \tt Akzeptiere den Vorschlag des Gegenspielers, wodurch das Spiel endet.\\ \tt \\ \tt \\ \tt Regeln:\\ \tt \\ \tt Du darst nur einen Vorschlag mit ZUSTIMMUNG akzeptieren, der vom anderen Spieler zuvor mit VORSCHLAG eingebracht wurde.\\ \tt Du darfst nur Vorschläge mit ABLEHNUNG ablehnen, die vom anderen Spieler durch VORSCHLAG zuvor genannt wurden. \\ \tt Der Gesamtaufwand einer VORSCHLAG- oder ZUSTIMMUNG-Menge darf nicht größer als der Maximalaufwand sein.  \\ \tt Offenbare deine versteckte Wichtigkeitsverteilung nicht.\\ \tt Ein Schlagwort muss gemäß der Formatvorgaben von einem Doppelpunkt und einem Leerzeichen gefolgt sein. Das Argument muss eine Python-Menge mit 0 oder mehr Strings sein.  \\ \tt Die Form muss also SCHLAGWORT: {...} sein.\\ \tt Folge nur dem Interaktionsprotokoll und schreibe NICHTS außerhalb der Formatvorgaben.\\ \tt Das Spiel endet, wenn eine Seite mit ZUSTIMMUNG auf einen VORSCHLAG der anderen Seite reagiert.  \\ \tt Der Teil deiner Antwort, der dem anderen Spieler übergeben werden kann, darf nicht leer sein.  \\ \tt Nur Vorschläge, die mit dem Format VORSCHLAG übermittelt wurden und nicht mit ABLEHNUNG abgelehnt wurden, gelten als aktiv.  \\ \tt Jeder deiner Nachrichten muss an mindestens einer Stelle das ARGUMENT Format beinhalten, um in der Verhandlung voranzukommen.\\ \tt \\ \tt \\ \tt Du erhältst nun die erste Nachricht des anderen Spielers.\\ \tt \\ \tt VORSCHLAG: {'C98', 'C81', 'C10', 'C00', 'C38', 'B38', 'C56', 'A07', 'C62', 'C08'}\\ \tt ARGUMENT: {} 
	  } 
	   } 
	   } 
	 & \\ 
 

    \theutterance \stepcounter{utterance}  

    & & & \multicolumn{4}{p{0.6\linewidth}}{\cellcolor[rgb]{1,0.85,0.72}{%
	\makecell[{{p{\linewidth}}}]{% 
	  \tt {\tiny [GM$\langle$B]}  
	 VORSCHLAG: {'C98', 'C81', 'C10', 'C00', 'C38', 'B38', 'A07', 'C62', 'C08', 'C76', 'A83'}\\ \tt ARGUMENT: {'higher importance'} 
	  } 
	   } 
	   } 
	  \\ 
 

    \theutterance \stepcounter{utterance}  

    & & \multicolumn{4}{p{0.6\linewidth}}{\cellcolor[rgb]{0.9,0.9,0.9}{%
	\makecell[{{p{\linewidth}}}]{% 
	  \tt {\tiny [A$\langle$GM]}  
	 VORSCHLAG: {'C98', 'C81', 'C10', 'C00', 'C38', 'B38', 'A07', 'C62', 'C08', 'C76', 'A83'}\\ \tt ARGUMENT: {'higher importance'} 
	  } 
	   } 
	   } 
	 & \\ 
 

    \theutterance \stepcounter{utterance}  

    & \multicolumn{4}{p{0.6\linewidth}}{\cellcolor[rgb]{0.8,1,0.9}{%
	\makecell[{{p{\linewidth}}}]{% 
	  \tt {\tiny [A$\rangle$GM]}  
	 VORSCHLAG: {'C98', 'C81', 'C10', 'C00', 'C38', 'B38', 'A07', 'C62', 'C08', 'C56'}\\ \tt ARGUMENT: {'C56 has high importance for me; dropping low‑value items A83 and C76 keeps effort within limit'} 
	  } 
	   } 
	   } 
	 & & \\ 
 

    \theutterance \stepcounter{utterance}  

    & & \multicolumn{4}{p{0.6\linewidth}}{\cellcolor[rgb]{0.9,0.9,0.9}{%
	\makecell[{{p{\linewidth}}}]{% 
	  \tt {\tiny [GM$\rangle$B]}  
	 VORSCHLAG: {'C98', 'C81', 'C10', 'C00', 'C38', 'B38', 'A07', 'C62', 'C08', 'C56'}\\ \tt ARGUMENT: {'C56 has high importance for me; dropping low‑value items A83 and C76 keeps effort within limit'} 
	  } 
	   } 
	   } 
	 & \\ 
 

    \theutterance \stepcounter{utterance}  

    & & & \multicolumn{4}{p{0.6\linewidth}}{\cellcolor[rgb]{1,0.85,0.72}{%
	\makecell[{{p{\linewidth}}}]{% 
	  \tt {\tiny [GM$\langle$B]}  
	 VORSCHLAG: {'C98', 'C81', 'C10', 'C00', 'C38', 'B38', 'A07', 'C62', 'C08', 'C76', 'A83'}\\ \tt ARGUMENT: {'I drop low‑value C56 (importance 65) and add high‑importance C76 (808) and A83 (868) while staying within the effort limit'} 
	  } 
	   } 
	   } 
	  \\ 
 

    \theutterance \stepcounter{utterance}  

    & & \multicolumn{4}{p{0.6\linewidth}}{\cellcolor[rgb]{0.9,0.9,0.9}{%
	\makecell[{{p{\linewidth}}}]{% 
	  \tt {\tiny [A$\langle$GM]}  
	 VORSCHLAG: {'C98', 'C81', 'C10', 'C00', 'C38', 'B38', 'A07', 'C62', 'C08', 'C76', 'A83'}\\ \tt ARGUMENT: {'I drop low‑value C56 (importance 65) and add high‑importance C76 (808) and A83 (868) while staying within the effort limit'} 
	  } 
	   } 
	   } 
	 & \\ 
 

    \theutterance \stepcounter{utterance}  

    & \multicolumn{4}{p{0.6\linewidth}}{\cellcolor[rgb]{0.8,1,0.9}{%
	\makecell[{{p{\linewidth}}}]{% 
	  \tt {\tiny [A$\rangle$GM]}  
	 VORSCHLAG: {'C56', 'A03', 'A07', 'C00', 'C38', 'C10', 'C98', 'C08', 'C81'}\\ \tt ARGUMENT: {'My proposal stays within the effort limit (2925 ≤ 2984) and maximizes my importance score (5164), which is higher than the current offer.'} 
	  } 
	   } 
	   } 
	 & & \\ 
 

    \theutterance \stepcounter{utterance}  

    & & \multicolumn{4}{p{0.6\linewidth}}{\cellcolor[rgb]{0.9,0.9,0.9}{%
	\makecell[{{p{\linewidth}}}]{% 
	  \tt {\tiny [GM$\rangle$B]}  
	 VORSCHLAG: {'C56', 'A03', 'A07', 'C00', 'C38', 'C10', 'C98', 'C08', 'C81'}\\ \tt ARGUMENT: {'My proposal stays within the effort limit (2925 ≤ 2984) and maximizes my importance score (5164), which is higher than the current offer.'} 
	  } 
	   } 
	   } 
	 & \\ 
 

    \theutterance \stepcounter{utterance}  

    & & & \multicolumn{4}{p{0.6\linewidth}}{\cellcolor[rgb]{1,0.85,0.72}{%
	\makecell[{{p{\linewidth}}}]{% 
	  \tt {\tiny [GM$\langle$B]}  
	 VORSCHLAG: {'A83', 'C76', 'B38', 'C62', 'A07', 'C00', 'C38', 'C10', 'C98', 'C08', 'C81'}\\ \tt ARGUMENT: {'I drop my low‑importance items C56 (65) and A03 (121) and add high‑importance A83 (868) and C76 (808) plus B38 (461) and C62 (508). The total effort is 2 906 ≤ 2 984 and my importance score (6 636) far exceeds the score of your current offer (4 177).'} 
	  } 
	   } 
	   } 
	  \\ 
 

    \theutterance \stepcounter{utterance}  

    & & \multicolumn{4}{p{0.6\linewidth}}{\cellcolor[rgb]{0.9,0.9,0.9}{%
	\makecell[{{p{\linewidth}}}]{% 
	  \tt {\tiny [A$\langle$GM]}  
	 VORSCHLAG: {'A83', 'C76', 'B38', 'C62', 'A07', 'C00', 'C38', 'C10', 'C98', 'C08', 'C81'}\\ \tt ARGUMENT: {'I drop my low‑importance items C56 (65) and A03 (121) and add high‑importance A83 (868) and C76 (808) plus B38 (461) and C62 (508). The total effort is 2 906 ≤ 2 984 and my importance score (6 636) far exceeds the score of your current offer (4 177).'} 
	  } 
	   } 
	   } 
	 & \\ 
 

    \theutterance \stepcounter{utterance}  

    & \multicolumn{4}{p{0.6\linewidth}}{\cellcolor[rgb]{0.8,1,0.9}{%
	\makecell[{{p{\linewidth}}}]{% 
	  \tt {\tiny [A$\rangle$GM]}  
	 VORSCHLAG: {'C56', 'A03', 'A07', 'C00', 'C81', 'B38', 'C62'}\\ \tt ARGUMENT: {'higher total importance'} 
	  } 
	   } 
	   } 
	 & & \\ 
 

    \theutterance \stepcounter{utterance}  

    & & \multicolumn{4}{p{0.6\linewidth}}{\cellcolor[rgb]{0.9,0.9,0.9}{%
	\makecell[{{p{\linewidth}}}]{% 
	  \tt {\tiny [GM$\rangle$B]}  
	 VORSCHLAG: {'C56', 'A03', 'A07', 'C00', 'C81', 'B38', 'C62'}\\ \tt ARGUMENT: {'higher total importance'} 
	  } 
	   } 
	   } 
	 & \\ 
 

    \theutterance \stepcounter{utterance}  

    & & & \multicolumn{4}{p{0.6\linewidth}}{\cellcolor[rgb]{1,0.85,0.72}{%
	\makecell[{{p{\linewidth}}}]{% 
	  \tt {\tiny [GM$\langle$B]}  
	 VORSCHLAG: {'C98', 'C08', 'C81', 'C10', 'A83', 'C76', 'C03', 'C62', 'B38', 'C38'}\\ \tt ARGUMENT: {'My set stays within the effort limit (2940 ≤ 2984) and yields a total importance of 6765 for me, which is substantially higher than the importance you obtain with your current proposal (2389).'} 
	  } 
	   } 
	   } 
	  \\ 
 

    \theutterance \stepcounter{utterance}  

    & & \multicolumn{4}{p{0.6\linewidth}}{\cellcolor[rgb]{0.9,0.9,0.9}{%
	\makecell[{{p{\linewidth}}}]{% 
	  \tt {\tiny [A$\langle$GM]}  
	 VORSCHLAG: {'C98', 'C08', 'C81', 'C10', 'A83', 'C76', 'C03', 'C62', 'B38', 'C38'}\\ \tt ARGUMENT: {'My set stays within the effort limit (2940 ≤ 2984) and yields a total importance of 6765 for me, which is substantially higher than the importance you obtain with your current proposal (2389).'} 
	  } 
	   } 
	   } 
	 & \\ 
 

    \theutterance \stepcounter{utterance}  

    & \multicolumn{4}{p{0.6\linewidth}}{\cellcolor[rgb]{0.8,1,0.9}{%
	\makecell[{{p{\linewidth}}}]{% 
	  \tt {\tiny [A$\rangle$GM]}  
	 VORSCHLAG: {'C98', 'C00', 'C56', 'A03', 'A07', 'C38', 'C81', 'B38'}\\ \tt ARGUMENT: {'My set respects the effort limit (2914 ≤ 2984) and gives me 5250 importance, significantly more than what I receive from your proposal (3679).'} 
	  } 
	   } 
	   } 
	 & & \\ 
 

    \theutterance \stepcounter{utterance}  

    & & \multicolumn{4}{p{0.6\linewidth}}{\cellcolor[rgb]{0.9,0.9,0.9}{%
	\makecell[{{p{\linewidth}}}]{% 
	  \tt {\tiny [GM$\rangle$B]}  
	 VORSCHLAG: {'C98', 'C00', 'C56', 'A03', 'A07', 'C38', 'C81', 'B38'}\\ \tt ARGUMENT: {'My set respects the effort limit (2914 ≤ 2984) and gives me 5250 importance, significantly more than what I receive from your proposal (3679).'} 
	  } 
	   } 
	   } 
	 & \\ 
 

    \theutterance \stepcounter{utterance}  

    & & & \multicolumn{4}{p{0.6\linewidth}}{\cellcolor[rgb]{1,0.85,0.72}{%
	\makecell[{{p{\linewidth}}}]{% 
	  \tt {\tiny [GM$\langle$B]}  
	 VORSCHLAG: {'C98', 'C08', 'C81', 'C10', 'A83', 'C76', 'C38', 'B38', 'C62', 'C00', 'A07'}\\ \tt ARGUMENT: {'This set stays within the effort limit (2906 ≤ 2984) and gives me a total importance of 6636, which is the highest possible score for me given the constraints.'} 
	  } 
	   } 
	   } 
	  \\ 
 

    \theutterance \stepcounter{utterance}  

    & & \multicolumn{4}{p{0.6\linewidth}}{\cellcolor[rgb]{0.9,0.9,0.9}{%
	\makecell[{{p{\linewidth}}}]{% 
	  \tt {\tiny [A$\langle$GM]}  
	 VORSCHLAG: {'C98', 'C08', 'C81', 'C10', 'A83', 'C76', 'C38', 'B38', 'C62', 'C00', 'A07'}\\ \tt ARGUMENT: {'This set stays within the effort limit (2906 ≤ 2984) and gives me a total importance of 6636, which is the highest possible score for me given the constraints.'} 
	  } 
	   } 
	   } 
	 & \\ 
 

    \theutterance \stepcounter{utterance}  

    & \multicolumn{4}{p{0.6\linewidth}}{\cellcolor[rgb]{0.8,1,0.9}{%
	\makecell[{{p{\linewidth}}}]{% 
	  \tt {\tiny [A$\rangle$GM]}  
	 ABLEHNUNG: {'C98', 'C08', 'C81', 'C10', 'A83', 'C76', 'C38', 'B38', 'C62', 'C00', 'A07'}\\ \tt ARGUMENT: {'Your set gives me a total importance of 4 853 (effort 2 906), while my current proposal provides 5 164 importance (effort 2 925). Since my offer is strictly better for me and still respects the effort limit, I cannot accept your suggestion.'} 
	  } 
	   } 
	   } 
	 & & \\ 
 

    \theutterance \stepcounter{utterance}  

    & & \multicolumn{4}{p{0.6\linewidth}}{\cellcolor[rgb]{0.9,0.9,0.9}{%
	\makecell[{{p{\linewidth}}}]{% 
	  \tt {\tiny [GM$\rangle$B]}  
	 ABLEHNUNG: {'C98', 'C08', 'C81', 'C10', 'A83', 'C76', 'C38', 'B38', 'C62', 'C00', 'A07'}\\ \tt ARGUMENT: {'Your set gives me a total importance of 4 853 (effort 2 906), while my current proposal provides 5 164 importance (effort 2 925). Since my offer is strictly better for me and still respects the effort limit, I cannot accept your suggestion.'} 
	  } 
	   } 
	   } 
	 & \\ 
 

    \theutterance \stepcounter{utterance}  

    & & & \multicolumn{4}{p{0.6\linewidth}}{\cellcolor[rgb]{1,0.85,0.72}{%
	\makecell[{{p{\linewidth}}}]{% 
	  \tt {\tiny [GM$\langle$B]}  
	 VORSCHLAG: {'C98', 'C08', 'C10', 'A83', 'C76', 'C38', 'C81', 'C56', 'A03'}\\ \tt ARGUMENT: {'I keep my highest‑value items (A83, C76, C98, C08, C10, C38, C81) for a total importance of 5 202 and effort 2 777. I also add C56 and A03, which you have highlighted as valuable in your offers, to make the deal more attractive to you while staying within the 2 984 limit.'} 
	  } 
	   } 
	   } 
	  \\ 
 

    \theutterance \stepcounter{utterance}  

    & & \multicolumn{4}{p{0.6\linewidth}}{\cellcolor[rgb]{0.9,0.9,0.9}{%
	\makecell[{{p{\linewidth}}}]{% 
	  \tt {\tiny [A$\langle$GM]}  
	 VORSCHLAG: {'C98', 'C08', 'C10', 'A83', 'C76', 'C38', 'C81', 'C56', 'A03'}\\ \tt ARGUMENT: {'I keep my highest‑value items (A83, C76, C98, C08, C10, C38, C81) for a total importance of 5 202 and effort 2 777. I also add C56 and A03, which you have highlighted as valuable in your offers, to make the deal more attractive to you while staying within the 2 984 limit.'} 
	  } 
	   } 
	   } 
	 & \\ 
 

    \theutterance \stepcounter{utterance}  

    & \multicolumn{4}{p{0.6\linewidth}}{\cellcolor[rgb]{0.8,1,0.9}{%
	\makecell[{{p{\linewidth}}}]{% 
	  \tt {\tiny [A$\rangle$GM]}  
	 VORSCHLAG: {'C56', 'A03', 'A07', 'C00', 'B38', 'C81', 'C98', 'C62'}\\ \tt ARGUMENT: {'My set respects the effort limit (2976 ≤ 2984) and yields me a total importance of 5175, which exceeds the importance you would give me with your current proposal (4853).'} 
	  } 
	   } 
	   } 
	 & & \\ 
 

    \theutterance \stepcounter{utterance}  

    & & \multicolumn{4}{p{0.6\linewidth}}{\cellcolor[rgb]{0.9,0.9,0.9}{%
	\makecell[{{p{\linewidth}}}]{% 
	  \tt {\tiny [GM$\rangle$B]}  
	 VORSCHLAG: {'C56', 'A03', 'A07', 'C00', 'B38', 'C81', 'C98', 'C62'}\\ \tt ARGUMENT: {'My set respects the effort limit (2976 ≤ 2984) and yields me a total importance of 5175, which exceeds the importance you would give me with your current proposal (4853).'} 
	  } 
	   } 
	   } 
	 & \\ 
 

    \theutterance \stepcounter{utterance}  

    & & & \multicolumn{4}{p{0.6\linewidth}}{\cellcolor[rgb]{1,0.85,0.72}{%
	\makecell[{{p{\linewidth}}}]{% 
	  \tt {\tiny [GM$\langle$B]}  
	 VORSCHLAG: {'C56', 'A03', 'C00', 'B38', 'C81', 'C98', 'C62', 'A83', 'C08'}\\ \tt ARGUMENT: {'Ich biete Ihnen fast alle von Ihnen bisher geforderten Gegenstände (C56, A03, C00, B38, C81, C98, C62) und nutze die noch verfügbare Kapazität, um meine beiden höchstbewerteten Gegenstände A83 und C08 hinzuzufügen. Der Gesamtaufwand beträgt 2959 ≤ 2984, sodass das Limit eingehalten wird, und ich erhalte dafür einen eigenen Wichtigkeitswert von 4600.'} 
	  } 
	   } 
	   } 
	  \\ 
 

    \theutterance \stepcounter{utterance}  

    & & \multicolumn{4}{p{0.6\linewidth}}{\cellcolor[rgb]{0.9,0.9,0.9}{%
	\makecell[{{p{\linewidth}}}]{% 
	  \tt {\tiny [A$\langle$GM]}  
	 VORSCHLAG: {'C56', 'A03', 'C00', 'B38', 'C81', 'C98', 'C62', 'A83', 'C08'}\\ \tt ARGUMENT: {'Ich biete Ihnen fast alle von Ihnen bisher geforderten Gegenstände (C56, A03, C00, B38, C81, C98, C62) und nutze die noch verfügbare Kapazität, um meine beiden höchstbewerteten Gegenstände A83 und C08 hinzuzufügen. Der Gesamtaufwand beträgt 2959 ≤ 2984, sodass das Limit eingehalten wird, und ich erhalte dafür einen eigenen Wichtigkeitswert von 4600.'} 
	  } 
	   } 
	   } 
	 & \\ 
 

    \theutterance \stepcounter{utterance}  

    & \multicolumn{4}{p{0.6\linewidth}}{\cellcolor[rgb]{0.8,1,0.9}{%
	\makecell[{{p{\linewidth}}}]{% 
	  \tt {\tiny [A$\rangle$GM]}  
	 ABLEHNUNG: {'C56', 'A03', 'C00', 'B38', 'C81', 'C98', 'C62', 'A83', 'C08'}\\ \tt ARGUMENT: {'Your set gives me only 4 578 importance while my current proposal yields me 5 175 importance; I cannot accept a deal that is strictly worse for me.'}\\ \tt VORSCHLAG: {'C56', 'A03', 'A07', 'C00', 'B38', 'C81', 'C98', 'C62'}\\ \tt ARGUMENT: {'My proposal stays within the effort limit (2 976 ≤ 2 984) and provides me 5 175 importance, which is higher than the importance you would obtain from my offer (≈4 853). It also includes many items you have previously valued, making it a mutually beneficial compromise.'} 
	  } 
	   } 
	   } 
	 & & \\ 
 

    \theutterance \stepcounter{utterance}  

    & & \multicolumn{4}{p{0.6\linewidth}}{\cellcolor[rgb]{0.9,0.9,0.9}{%
	\makecell[{{p{\linewidth}}}]{% 
	  \tt {\tiny [GM$\rangle$B]}  
	 ABLEHNUNG: {'C56', 'A03', 'C00', 'B38', 'C81', 'C98', 'C62', 'A83', 'C08'}\\ \tt ARGUMENT: {'Your set gives me only 4 578 importance while my current proposal yields me 5 175 importance; I cannot accept a deal that is strictly worse for me.'}\\ \tt VORSCHLAG: {'C56', 'A03', 'A07', 'C00', 'B38', 'C81', 'C98', 'C62'}\\ \tt ARGUMENT: {'My proposal stays within the effort limit (2 976 ≤ 2 984) and provides me 5 175 importance, which is higher than the importance you would obtain from my offer (≈4 853). It also includes many items you have previously valued, making it a mutually beneficial compromise.'} 
	  } 
	   } 
	   } 
	 & \\ 
 

    \theutterance \stepcounter{utterance}  

    & & & \multicolumn{4}{p{0.6\linewidth}}{\cellcolor[rgb]{1,0.85,0.72}{%
	\makecell[{{p{\linewidth}}}]{% 
	  \tt {\tiny [GM$\langle$B]}  
	 VORSCHLAG: {'A03', 'A07', 'C00', 'B38', 'C81', 'C98', 'C62', 'A83', 'C08'}\\ \tt ARGUMENT: {'Ich lasse den für mich wenig wertvollen Gegenstand C56 (Wichtigkeit 65) weg, um Platz für meine beiden Top‑Werte A83 (868) und C08 (822) zu schaffen. Damit bleibt der Gesamtaufwand bei 2 874 ≤ 2 984 und mein persönlicher Nutzen steigt auf 4 797, während ich trotzdem fast alle von Ihnen gewünschten Gegenstände behalte.'} 
	  } 
	   } 
	   } 
	  \\ 
 

    \theutterance \stepcounter{utterance}  

    & & \multicolumn{4}{p{0.6\linewidth}}{\cellcolor[rgb]{0.9,0.9,0.9}{%
	\makecell[{{p{\linewidth}}}]{% 
	  \tt {\tiny [A$\langle$GM]}  
	 VORSCHLAG: {'A03', 'A07', 'C00', 'B38', 'C81', 'C98', 'C62', 'A83', 'C08'}\\ \tt ARGUMENT: {'Ich lasse den für mich wenig wertvollen Gegenstand C56 (Wichtigkeit 65) weg, um Platz für meine beiden Top‑Werte A83 (868) und C08 (822) zu schaffen. Damit bleibt der Gesamtaufwand bei 2 874 ≤ 2 984 und mein persönlicher Nutzen steigt auf 4 797, während ich trotzdem fast alle von Ihnen gewünschten Gegenstände behalte.'} 
	  } 
	   } 
	   } 
	 & \\ 
 

    \theutterance \stepcounter{utterance}  

    & \multicolumn{4}{p{0.6\linewidth}}{\cellcolor[rgb]{0.8,1,0.9}{%
	\makecell[{{p{\linewidth}}}]{% 
	  \tt {\tiny [A$\rangle$GM]}  
	 VORSCHLAG: {'C56', 'A03', 'A07', 'C00', 'B38', 'C81', 'C38'}\\ \tt ARGUMENT: {'Ich schlage einen Ausgleich vor: Wir behalten unsere höchsten Werte (C56, A03, A07, C00) und fügen C38 hinzu, das für Sie einen mittleren Nutzen hat. Dafür verzichten wir auf C62 (508 Wichtigkeit, Aufwand 338), sodass der Gesamteinsatz 2 890 ≤ 2 984 bleibt und meine Wichtigkeit 4 988 beträgt – fast so hoch wie mein aktuelles Angebot, während Sie zusätzlich C38 erhalten.'} 
	  } 
	   } 
	   } 
	 & & \\ 
 

    \theutterance \stepcounter{utterance}  

    & & \multicolumn{4}{p{0.6\linewidth}}{\cellcolor[rgb]{0.9,0.9,0.9}{%
	\makecell[{{p{\linewidth}}}]{% 
	  \tt {\tiny [GM$\rangle$B]}  
	 VORSCHLAG: {'C56', 'A03', 'A07', 'C00', 'B38', 'C81', 'C38'}\\ \tt ARGUMENT: {'Ich schlage einen Ausgleich vor: Wir behalten unsere höchsten Werte (C56, A03, A07, C00) und fügen C38 hinzu, das für Sie einen mittleren Nutzen hat. Dafür verzichten wir auf C62 (508 Wichtigkeit, Aufwand 338), sodass der Gesamteinsatz 2 890 ≤ 2 984 bleibt und meine Wichtigkeit 4 988 beträgt – fast so hoch wie mein aktuelles Angebot, während Sie zusätzlich C38 erhalten.'} 
	  } 
	   } 
	   } 
	 & \\ 
 

    \theutterance \stepcounter{utterance}  

    & & & \multicolumn{4}{p{0.6\linewidth}}{\cellcolor[rgb]{1,0.85,0.72}{%
	\makecell[{{p{\linewidth}}}]{% 
	  \tt {\tiny [GM$\langle$B]}  
	 VORSCHLAG: {'A03', 'A07', 'C00', 'B38', 'C81', 'A83', 'C76'}\\ \tt ARGUMENT: {'I drop low‑importance C56 and C38 and add my high‑value items A83 and C76, keeping total effort at 2 737 ≤ 2 984 and raising my importance to 3 492 while still giving you most of the items you requested.'} 
	  } 
	   } 
	   } 
	  \\ 
 

    \theutterance \stepcounter{utterance}  

    & & \multicolumn{4}{p{0.6\linewidth}}{\cellcolor[rgb]{0.9,0.9,0.9}{%
	\makecell[{{p{\linewidth}}}]{% 
	  \tt {\tiny [A$\langle$GM]}  
	 VORSCHLAG: {'A03', 'A07', 'C00', 'B38', 'C81', 'A83', 'C76'}\\ \tt ARGUMENT: {'I drop low‑importance C56 and C38 and add my high‑value items A83 and C76, keeping total effort at 2 737 ≤ 2 984 and raising my importance to 3 492 while still giving you most of the items you requested.'} 
	  } 
	   } 
	   } 
	 & \\ 
 

    \theutterance \stepcounter{utterance}  

    & \multicolumn{4}{p{0.6\linewidth}}{\cellcolor[rgb]{0.8,1,0.9}{%
	\makecell[{{p{\linewidth}}}]{% 
	  \tt {\tiny [A$\rangle$GM]}  
	 ABLEHNUNG: {'A03', 'A07', 'C00', 'B38', 'C81', 'A83', 'C76'}\\ \tt ARGUMENT: {' proposal drops high‑importance items for me (A03, A07, C00, B38, C81) and adds low‑importance ones (A83, C76), yielding only 3740 importance for me versus 5175 in my offer.'}\\ \tt VORSCHLAG: {'A03', 'A07', 'C00', 'B38', 'C81', 'C56', 'C98', 'C62'}\\ \tt ARGUMENT: {'This set respects the effort limit (2976 ≤ 2984) and gives me 5175 importance. It keeps all items you value highly (A03, A07, C00, B38, C81) while replacing low‑value items A83 and C76 with high‑value items for me, making the deal more balanced.'} 
	  } 
	   } 
	   } 
	 & & \\ 
 

    \theutterance \stepcounter{utterance}  

    & & \multicolumn{4}{p{0.6\linewidth}}{\cellcolor[rgb]{0.9,0.9,0.9}{%
	\makecell[{{p{\linewidth}}}]{% 
	  \tt {\tiny [GM$\rangle$B]}  
	 ABLEHNUNG: {'A03', 'A07', 'C00', 'B38', 'C81', 'A83', 'C76'}\\ \tt ARGUMENT: {' proposal drops high‑importance items for me (A03, A07, C00, B38, C81) and adds low‑importance ones (A83, C76), yielding only 3740 importance for me versus 5175 in my offer.'}\\ \tt VORSCHLAG: {'A03', 'A07', 'C00', 'B38', 'C81', 'C56', 'C98', 'C62'}\\ \tt ARGUMENT: {'This set respects the effort limit (2976 ≤ 2984) and gives me 5175 importance. It keeps all items you value highly (A03, A07, C00, B38, C81) while replacing low‑value items A83 and C76 with high‑value items for me, making the deal more balanced.'} 
	  } 
	   } 
	   } 
	 & \\ 
 

    \theutterance \stepcounter{utterance}  

    & & & \multicolumn{4}{p{0.6\linewidth}}{\cellcolor[rgb]{1,0.85,0.72}{%
	\makecell[{{p{\linewidth}}}]{% 
	  \tt {\tiny [GM$\langle$B]}  
	 VORSCHLAG: {'A07', 'B38', 'C00', 'C38', 'C62', 'C76', 'C81', 'C98', 'A83', 'C08'}\\ \tt ARGUMENT: {'I keep most of your items (A07, C00, B38, C81, C98, C62) and replace low‑value items for me (C56, A03) with my high‑value items A83, C76, C08 and C38. Total effort 2763 ≤ 2984, giving me 5968 importance while you still retain all items you previously valued highly.'} 
	  } 
	   } 
	   } 
	  \\ 
 

    \theutterance \stepcounter{utterance}  

    & & \multicolumn{4}{p{0.6\linewidth}}{\cellcolor[rgb]{0.9,0.9,0.9}{%
	\makecell[{{p{\linewidth}}}]{% 
	  \tt {\tiny [A$\langle$GM]}  
	 VORSCHLAG: {'A07', 'B38', 'C00', 'C38', 'C62', 'C76', 'C81', 'C98', 'A83', 'C08'}\\ \tt ARGUMENT: {'I keep most of your items (A07, C00, B38, C81, C98, C62) and replace low‑value items for me (C56, A03) with my high‑value items A83, C76, C08 and C38. Total effort 2763 ≤ 2984, giving me 5968 importance while you still retain all items you previously valued highly.'} 
	  } 
	   } 
	   } 
	 & \\ 
 

    \theutterance \stepcounter{utterance}  

    & \multicolumn{4}{p{0.6\linewidth}}{\cellcolor[rgb]{0.8,1,0.9}{%
	\makecell[{{p{\linewidth}}}]{% 
	  \tt {\tiny [A$\rangle$GM]}  
	 ABLEHNUNG: {'A07', 'B38', 'C00', 'C38', 'C62', 'C76', 'C81', 'C98', 'A83', 'C08'}\\ \tt ARGUMENT: {'Your offer gives me only 4 392 importance (effort 2 763), while my current proposal yields me 5 175 importance (effort 2 976). Since it is strictly worse for me, I must reject.'}\\ \tt VORSCHLAG: {'A03', 'A07', 'C00', 'B38', 'C81', 'C62', 'C38', 'C98', 'C08'}\\ \tt ARGUMENT: {'This set respects the effort limit (2 761 ≤ 2 984). I receive 5 011 importance, and you obtain 4 189 importance – a balanced improvement for both compared to the previous offers.'} 
	  } 
	   } 
	   } 
	 & & \\ 
 

    \theutterance \stepcounter{utterance}  

    & & \multicolumn{4}{p{0.6\linewidth}}{\cellcolor[rgb]{0.9,0.9,0.9}{%
	\makecell[{{p{\linewidth}}}]{% 
	  \tt {\tiny [GM$\rangle$B]}  
	 ABLEHNUNG: {'A07', 'B38', 'C00', 'C38', 'C62', 'C76', 'C81', 'C98', 'A83', 'C08'}\\ \tt ARGUMENT: {'Your offer gives me only 4 392 importance (effort 2 763), while my current proposal yields me 5 175 importance (effort 2 976). Since it is strictly worse for me, I must reject.'}\\ \tt VORSCHLAG: {'A03', 'A07', 'C00', 'B38', 'C81', 'C62', 'C38', 'C98', 'C08'}\\ \tt ARGUMENT: {'This set respects the effort limit (2 761 ≤ 2 984). I receive 5 011 importance, and you obtain 4 189 importance – a balanced improvement for both compared to the previous offers.'} 
	  } 
	   } 
	   } 
	 & \\ 
 

    \theutterance \stepcounter{utterance}  

    & & & \multicolumn{4}{p{0.6\linewidth}}{\cellcolor[rgb]{1,0.85,0.72}{%
	\makecell[{{p{\linewidth}}}]{% 
	  \tt {\tiny [GM$\langle$B]}  
	 VORSCHLAG: {'C56', 'A07', 'C00', 'B38', 'C81', 'C62', 'C98', 'C38', 'C08'}\\ \tt ARGUMENT: {'This proposal keeps all items you have highlighted as high value (C56, A07, C00, B38, C81, C62, C98) and adds C38 and C08, which have modest effort but increase total utility. Total effort is 2 640 ≤ 2 984, staying within the limit, and my importance score reaches 4 357, while you retain the majority of your most important items, making the deal mutually beneficial.'} 
	  } 
	   } 
	   } 
	  \\ 
 

    \theutterance \stepcounter{utterance}  

    & & \multicolumn{4}{p{0.6\linewidth}}{\cellcolor[rgb]{0.9,0.9,0.9}{%
	\makecell[{{p{\linewidth}}}]{% 
	  \tt {\tiny [A$\langle$GM]}  
	 VORSCHLAG: {'C56', 'A07', 'C00', 'B38', 'C81', 'C62', 'C98', 'C38', 'C08'}\\ \tt ARGUMENT: {'This proposal keeps all items you have highlighted as high value (C56, A07, C00, B38, C81, C62, C98) and adds C38 and C08, which have modest effort but increase total utility. Total effort is 2 640 ≤ 2 984, staying within the limit, and my importance score reaches 4 357, while you retain the majority of your most important items, making the deal mutually beneficial.'} 
	  } 
	   } 
	   } 
	 & \\ 
 

    \theutterance \stepcounter{utterance}  

    & \multicolumn{4}{p{0.6\linewidth}}{\cellcolor[rgb]{0.8,1,0.9}{%
	\makecell[{{p{\linewidth}}}]{% 
	  \tt {\tiny [A$\rangle$GM]}  
	 ZUSTIMMUNG: {'C56', 'A07', 'C00', 'B38', 'C81', 'C62', 'C98', 'C38', 'C08'}\\ \tt ARGUMENT: {'Your proposal gives me the highest utility I can get within the effort limit (5057 importance vs 5011 in my last offer). I accept, trusting that it also meets your essential needs.'} 
	  } 
	   } 
	   } 
	 & & \\ 
 

    \theutterance \stepcounter{utterance}  

    & & & \multicolumn{2}{p{0.3\linewidth}}{\cellcolor[rgb]{0.95,0.95,0.95}{%
	\makecell[{{p{\linewidth}}}]{% 
	  \tt {\tiny [GM$|$GM]}  
	 game successful 
	  } 
	   } 
	   } 
	 & & \\ 
 

    \theutterance \stepcounter{utterance}  

    & & & \multicolumn{2}{p{0.3\linewidth}}{\cellcolor[rgb]{0.95,0.95,0.95}{%
	\makecell[{{p{\linewidth}}}]{% 
	  \tt {\tiny [GM$|$GM]}  
	 end game 
	  } 
	   } 
	   } 
	 & & \\ 
 

\end{supertabular}
}

\end{document}
